% =============================================================================
%  MCMOT Methods Report — Intra-camera Trackers & Inter-camera Association
% =============================================================================
\documentclass[11pt,a4paper]{article}

\usepackage[margin=2.2cm]{geometry}
\usepackage{amsmath,amssymb,mathtools}
\usepackage{booktabs,multirow,array,tabularx}
\usepackage{xcolor}
\usepackage{hyperref}
\usepackage{enumitem}
\usepackage{titlesec}
\usepackage{parskip}
\usepackage{microtype}
\usepackage{algorithm}
\usepackage{algpseudocode}
\usepackage{graphicx}

% ---------- colours -----------------------------------------------------------
\definecolor{navyblue}{HTML}{1a2744}
\definecolor{midblue}{HTML}{2563eb}
\definecolor{lightgray}{HTML}{f1f5f9}
\definecolor{darkgray}{HTML}{475569}
\definecolor{bestgreen}{HTML}{dcfce7}
\definecolor{alertred}{HTML}{fef2f2}

% ---------- section styling ---------------------------------------------------
\titleformat{\section}
  {\large\bfseries\color{navyblue}}
  {\thesection}{0.8em}{}
  [\color{navyblue}\titlerule\vspace{2pt}]

\titleformat{\subsection}
  {\normalsize\bfseries\color{midblue}}{\thesubsection}{0.6em}{}

% Starred subsubsections only — no counter, plain dark gray bold heading.
\titleformat{\subsubsection}
  {\normalsize\bfseries\color{darkgray}}{}{0pt}{}

% ---------- misc macros -------------------------------------------------------
\newcommand{\code}[1]{\texttt{\small#1}}
\newcommand{\lit}[1]{\textbf{\small[#1]}}
\newcommand{\Inf}{\infty}
\newcolumntype{L}[1]{>{\raggedright\arraybackslash}p{#1}}
\newcolumntype{C}[1]{>{\centering\arraybackslash}p{#1}}

% =============================================================================
\begin{document}

% ---------- Title block -------------------------------------------------------
\begin{center}
  {\LARGE\bfseries\color{navyblue}
    Multi-Camera Multi-Object Tracking\\[4pt]
    Methods \& Implementation Reference}\\[8pt]
  {\large Polybag Detection \& Tracking — OVGU AMS}\\[4pt]
  {\color{darkgray} July 2026}
\end{center}

\vspace{4pt}
\noindent\rule{\linewidth}{1.5pt}
\vspace{4pt}

\noindent
This document describes every tracking and inter-camera association method
used in the MCMOT pipeline, from the mathematical formulation to the
concrete implementation.  Methods are organised into three layers:
(1)~intra-camera object tracking,
(2)~offline inter-camera association, and
(3)~online (real-time) inter-camera association.

% =============================================================================
\section{Method Taxonomy}

\begin{table}[h]
\centering
\renewcommand{\arraystretch}{1.35}
\begin{tabularx}{\linewidth}{L{2.8cm} L{2.8cm} C{1.5cm} C{1.5cm} L{5.2cm}}
\toprule
\textbf{Method} & \textbf{Stage} & \textbf{Online?} & \textbf{Memory} & \textbf{Literature basis} \\
\midrule
ByteTrack  & Intra-camera  & Yes & $O(N)$ & Zhang et al.\ ECCV 2022 \lit{1} \\
BoT-SORT   & Intra-camera  & Yes & $O(N)$ & Aharon et al.\ arXiv 2022 \lit{2} \\
\midrule
\code{no\_assoc}          & Offline MCMOT & —  & none   & Baseline (identity) \\
\code{class\_only}        & Offline MCMOT & —  & $O(TN)$ & POM-inspired \lit{3} \\
\code{trk\_temporal}      & Offline MCMOT & —  & $O(TN)$ & JPDAF-inspired \lit{4} \\
\code{trk\_spatial}       & Offline MCMOT & —  & $O(TN)$ & POM-inspired \lit{3} \\
\code{trk\_combined}      & Offline MCMOT & —  & $O(TN)$ & LMGP/DyGLIP-inspired \lit{5} \\
\midrule
\code{class\_rank}        & Online MCMOT  & Yes & $O(N)$ & SORT variant \lit{6} \\
\code{class\_iou}         & Online MCMOT  & Yes & $O(N)$ & Hungarian / LAPJV \lit{7} \\
\code{class\_smooth}      & Online MCMOT  & Yes & $O(NW)$ & EMA / sliding-window vote \\
\bottomrule
\end{tabularx}
\caption{Complete taxonomy of all methods.  $N$ = number of active tracks,
  $T$ = sequence length, $W$ = smoothing window (8 frames).}
\end{table}

% =============================================================================
\section{Intra-Camera Trackers (Stage 1)}

Both trackers follow the \emph{Tracking-by-Detection} (TbD) paradigm:
detections from a YOLO OBB model are associated frame-to-frame using a
Kalman filter for motion prediction and the Hungarian algorithm for
minimum-cost assignment.

% -----------------------------------------------------------------------------
\subsection{ByteTrack}

\textbf{Source:} Zhang et al., ``ByteTrack: Multi-Object Tracking by Associating
Every Detection Box'', ECCV 2022.

\subsubsection*{Key idea}
Classical trackers discard low-confidence detections before association.
ByteTrack retains \emph{every} detection box and performs a two-stage match:
high-confidence detections first, then low-confidence ones.  This recovers
occluded or briefly-invisible objects that would otherwise be lost.

\subsubsection*{State representation}
Each track maintains a Kalman filter with state vector
\begin{equation}
  \mathbf{x}_k =
  \bigl[u,\; v,\; s,\; r,\; \dot{u},\; \dot{v},\; \dot{s}\bigr]^\top
  \label{eq:kalman-state}
\end{equation}
where $(u, v)$ is the bounding-box centre, $s$ is the area, $r$ is the
fixed aspect ratio, and $(\dot{u}, \dot{v}, \dot{s})$ are the corresponding
velocities.  The constant-velocity state-transition matrix is
\begin{equation}
  F = \begin{bmatrix} I_4 & \Delta t \cdot I_4 \\ 0 & I_4 \end{bmatrix},
  \qquad
  H = \begin{bmatrix} I_4 & 0 \end{bmatrix}.
\end{equation}
Prediction step: $\hat{\mathbf{x}}_k = F\,\mathbf{x}_{k-1}$,\;
$\hat{P}_k = F P_{k-1} F^\top + Q$.
Update step: standard Kalman gain applied when a detection is matched.

\subsubsection*{Association}
Let $\mathcal{D}^H = \{d : \text{conf}(d) \ge \tau_H\}$ (high-confidence)
and $\mathcal{D}^L = \{d : \tau_L \le \text{conf}(d) < \tau_H\}$
(low-confidence) with $\tau_H = 0.6$, $\tau_L = 0.1$.

The \textbf{IoU cost} between prediction $\hat{b}_i$ and detection $b_j$ is
\begin{equation}
  C^{\text{IoU}}_{ij} = 1 - \text{IoU}\bigl(\hat{b}_i,\, b_j\bigr).
\end{equation}
\textbf{Stage 1}: Hungarian assignment on $\mathcal{T}_\text{active} \times
\mathcal{D}^H$ with threshold $C^{\text{IoU}} < 0.5$.\\
\textbf{Stage 2}: Unmatched tracks from Stage 1 are matched against
$\mathcal{D}^L$ using the same cost.\\
Tracks unmatched for $\delta$ frames are removed; detections unmatched in
both stages initialise new tracks.

\begin{algorithm}[H]
\caption{ByteTrack per-frame update}
\begin{algorithmic}[1]
\Require detections $\mathcal{D}_t$, active tracks $\mathcal{T}$
\State $\mathcal{D}^H \gets \{d \in \mathcal{D}_t : \text{conf}(d) \ge \tau_H\}$
\State $\mathcal{D}^L \gets \{d \in \mathcal{D}_t : \tau_L \le \text{conf}(d) < \tau_H\}$
\State Predict $\hat{\mathbf{x}}_i$ for each $i \in \mathcal{T}$ via Kalman
\State $M_1 \gets \textsc{Hungarian}(\mathcal{T}, \mathcal{D}^H, C^{\text{IoU}}, 0.5)$
\State $\mathcal{T}_u \gets \mathcal{T} \setminus \{\text{matched in } M_1\}$
\State $M_2 \gets \textsc{Hungarian}(\mathcal{T}_u, \mathcal{D}^L, C^{\text{IoU}}, 0.5)$
\State Update Kalman states for all matched tracks
\State Spawn new tracks for unmatched detections in $\mathcal{D}^H$
\State Remove tracks lost for $> \delta$ frames
\end{algorithmic}
\end{algorithm}

% -----------------------------------------------------------------------------
\subsection{BoT-SORT}

\textbf{Source:} Aharon et al., ``BoT-SORT: Robust Associations Multi-Pedestrian
Tracking'', arXiv:2206.14651, 2022.

\subsubsection*{Key idea}
BoT-SORT extends ByteTrack with two additions: (1)~\textbf{Global Motion
Compensation (GMC)}, which corrects Kalman predictions for camera ego-motion
using a homography estimate; and (2)~optional \textbf{Re-ID features} that
provide appearance-based matching alongside IoU.

\subsubsection*{Camera motion compensation}
Let $H_t \in \mathbb{R}^{3\times 3}$ be the homography between frames $t-1$
and $t$ (estimated via sparse optical flow / ECC).  Before association the
predicted bounding-box centre is corrected:
\begin{equation}
  \tilde{u}_i = \bigl[H_t \cdot (u_i,\, v_i,\, 1)^\top\bigr]_{1:2}.
\end{equation}
This makes the IoU cost robust to global scene motion.

\subsubsection*{Fused cost (with Re-ID)}
When appearance embeddings $\mathbf{e}_i, \mathbf{e}_j$ are available:
\begin{equation}
  C^{\text{fused}}_{ij}
  = \lambda_1\,C^{\text{IoU}}_{ij}
  + (1-\lambda_1)\,\bigl(1 - \cos(\mathbf{e}_i,\,\mathbf{e}_j)\bigr),
  \qquad \lambda_1 \in [0,1].
\end{equation}
In our setup Re-ID is disabled ($\lambda_1 = 1$); BoT-SORT therefore differs
from ByteTrack primarily through the GMC correction.

% =============================================================================
\section{Offline Inter-Camera Association (Stage 2)}

After intra-camera tracking, each camera $c \in \{1,2,3,4\}$ produces a set
of \emph{local tracklets}
$\mathcal{T}_c = \{T_{c,\ell}\}$, where each tracklet carries:
a local ID $\ell$, a class label $\text{cls}(T)$, a set of active frames
$F_T \subseteq \{1,\ldots,K\}$, and per-frame axis-aligned bounding boxes.
The goal is to assign a consistent \emph{global ID} $g(T)$ such that the
same physical object receives the same ID in all cameras.

% -----------------------------------------------------------------------------
\subsection{\texttt{no\_assoc} — Baseline}

No inter-camera step is performed.  Local IDs are promoted to global IDs with
a per-camera offset to avoid collisions:
\begin{equation}
  g(T_{c,\ell}) = \ell + (c-1) \times 1000.
\end{equation}
This gives an IDF1 near chance level ($\approx 25\%$) and serves as the
lower-bound baseline.

% -----------------------------------------------------------------------------
\subsection{\texttt{class\_only} — Frame-Level Rank Matching}

\textbf{Closest literature:} POM (Probabilistic Occupancy Map,
Fleuret et al.\ 2008 \lit{3}) — exploits a shared scene coordinate without
requiring calibration; here the shared coordinate is the x-axis rank.

\textbf{For each frame $t$ and class $c$}, collect detections
$\mathcal{D}^c_{t,i} = \{(x^{(k)}_c, \ell^{(k)})\}$ per camera $i$,
sorted by x-centre.  Assign the same global ID to equal-rank detections:
\begin{equation}
  g\!\left(T_{i,\,\ell^{(r)}_i}\right) = g\!\left(T_{j,\,\ell^{(r)}_j}\right)
  \quad \forall\; i \ne j, \quad r = 0,1,\ldots
\end{equation}
Global IDs persist across frames via a \code{local\_to\_global} dictionary;
a new ID is minted only when a (camera, local\_id) pair has not been seen before.

% -----------------------------------------------------------------------------
\subsection{\texttt{trk\_temporal} — Tracklet Temporal Jaccard}

\textbf{Closest literature:} JPDAF (Joint Probabilistic Data Association
Filter, Bar-Shalom \& Fortmann 1988 \lit{4}), which models the probability
that two detections belong to the same target via their temporal overlap.

\subsubsection*{Temporal Jaccard distance}
\begin{equation}
  d_J(T_a, T_b) = 1 - J(T_a, T_b), \qquad
  J(T_a, T_b) = \frac{|F_a \cap F_b|}{|F_a \cup F_b|}.
  \label{eq:jaccard}
\end{equation}
Because all cameras are synchronised and observe a shared scene, two tracklets
from different cameras that co-exist in many frames are likely the same object.

\subsubsection*{Assignment}
For each color class $c$, build a cost matrix
\begin{equation}
  C_{ab} =
  \begin{cases}
    d_J(T_a, T_b) & \text{if } \text{cls}(T_a) = \text{cls}(T_b)
                    \text{ and } \text{cam}(T_a) \ne \text{cam}(T_b)\\
    \infty & \text{otherwise}
  \end{cases}
\end{equation}
and apply greedy minimum-cost merging (pairs sorted ascending by $C_{ab}$),
enforcing the constraint that each camera contributes at most one tracklet
per global group.

% -----------------------------------------------------------------------------
\subsection{\texttt{trk\_spatial} — Tracklet Spatial Rank}

\textbf{Closest literature:} POM-style occupancy reasoning \lit{3}; also
related to the spatial ordering prior in StrongSORT/OC-SORT.

\subsubsection*{Per-camera mean x-centre}
\begin{equation}
  \bar{x}(T) = \frac{1}{|F_T|}
  \sum_{t \in F_T} \!\left(x_t + \tfrac{w_t}{2}\right).
\end{equation}
Within class $c$ and camera $i$, sort tracklets by $\bar{x}$ and assign a
\textit{normalised rank}:
\begin{equation}
  \hat{x}(T_{i,k}) = \frac{\text{rank}(T_{i,k})}{\max(|\mathcal{T}^c_i|-1,\,1)}.
\end{equation}
The spatial cost is then
\begin{equation}
  d_S(T_a, T_b) = \bigl|\hat{x}(T_a) - \hat{x}(T_b)\bigr| \in [0,1].
\end{equation}
The intuition is that objects on a conveyor belt preserve their left-to-right
order across all camera views, so equal normalised rank implies same object.

% -----------------------------------------------------------------------------
\subsection{\texttt{trk\_combined} — Combined Cost}

\textbf{Closest literature:} LMGP (Local Motion and Global Perception,
Tang et al.\ 2017 \lit{5}) and DyGLIP, which fuse appearance, motion, and
spatial cues without deep Re-ID features.

A simple equal-weight fusion of the two complementary cues:
\begin{equation}
  \boxed{d_{\text{comb}}(T_a, T_b) = 0.5\,d_J(T_a, T_b) + 0.5\,d_S(T_a, T_b).}
  \label{eq:combined}
\end{equation}
Greedy Hungarian merging is applied as in the individual methods.

% =============================================================================
\section{Online (Real-Time) Inter-Camera Association (Stage 2)}

The online pipeline processes all four cameras in parallel threads at each
timestep $t$ (implemented via \code{ThreadPoolExecutor}).  Association runs
on the \emph{currently active} track set, with no future information.

\vspace{4pt}
\noindent
Each camera returns a list of active tracks at time $t$:
\begin{equation*}
  \mathcal{A}_t^{(c)} = \left\{(\ell,\; \text{cls},\; x_c,\; \text{corners})\right\}.
\end{equation*}
A global ID mapping $\phi: (\text{cam}, \text{local\_id}) \to \text{global\_id}$
is maintained across frames and updated each timestep.

% -----------------------------------------------------------------------------
\subsection{\texttt{class\_rank} — Per-Frame x-Rank}

\textbf{Closest literature:} Simplified SORT (Bewley et al.\ 2016 \lit{6})
without the Kalman prediction, relying solely on the spatial ordering
constraint of a conveyor-belt scene.

\textbf{For each frame $t$ and class $c$}:
\begin{enumerate}[leftmargin=*,itemsep=2pt]
  \item For each camera $i$, collect active tracks of class $c$ and sort by
        x-centre:
        \begin{equation}
          \Pi^c_i = \text{argsort}\!\left(\{x^{(k)}_{c,i}\}_{k}\right).
        \end{equation}
  \item Assign the same global ID to same-rank tracks across cameras:
        \begin{equation}
          \phi\!\left(i,\, \Pi^c_i[r]\right)
          = \phi\!\left(j,\, \Pi^c_j[r]\right)
          \quad \forall\; i \ne j,\; r = 0, 1, \ldots
        \end{equation}
  \item If a (cam, local\_id) pair has no existing global ID, a new one is
        minted.
\end{enumerate}

\textbf{Complexity}: $O(N \log N)$ per class per frame.  Zero additional
memory beyond $\phi$.  Works perfectly when same-class objects never swap
their left-to-right order between cameras.

% -----------------------------------------------------------------------------
\subsection{\texttt{class\_iou} — Spatial Hungarian (Normalised $\Delta x$)}

\textbf{Closest literature:} Hungarian algorithm \lit{7} applied to a
one-dimensional normalised position difference.
\textit{Note}: the name ``class\_iou'' is a misnomer from an earlier design;
pixel-space IoU between cameras is always $\approx 0$ for uncalibrated views.
The actual cost is normalised $|\Delta x|$.

\subsubsection*{Per-frame assignment}
Let $\tilde{x}_{c,\ell} = x_{c,\ell} / W$ be the \emph{normalised} x-centre
(where $W = 1920$ px).  Split active class-$c$ tracks into two groups:
\begin{align}
  \mathcal{K} &= \{(c, \ell) : \phi(c,\ell) \text{ is already known}\}, \\
  \mathcal{U} &= \{(c, \ell) : \phi(c,\ell) \text{ is unknown}\}.
\end{align}
Build cost matrix $C \in \mathbb{R}^{|\mathcal{U}| \times |\mathcal{K}|}$:
\begin{equation}
  C_{uk} =
  \begin{cases}
    \bigl|\tilde{x}_u - \tilde{x}_k\bigr|
      & \text{if } \text{cam}(u) \ne \text{cam}(k) \\
    1 & \text{(blocked: same camera)}
  \end{cases}
\end{equation}
Solve the assignment problem:
\begin{equation}
  \sigma^* = \arg\min_{\sigma} \sum_{u \in \mathcal{U}} C_{u,\,\sigma(u)},
\end{equation}
accepting matches where $C_{u,\sigma^*(u)} < \tau = 0.30$
(i.e., within $\approx 576$ pixels at full resolution).  Unmatched unknowns
receive new global IDs.  When $\mathcal{K} = \emptyset$ (first frame), falls
back to \texttt{class\_rank}.

\textbf{Key difference from \texttt{class\_rank}}: rank matching is a strict
ordinal operation; spatial Hungarian \emph{minimises total displacement},
tolerating small rank inversions when two bags are spatially close.

% -----------------------------------------------------------------------------
\subsection{\texttt{class\_smooth} — Class Rank with Temporal Smoothing}

\textbf{Motivation:} \texttt{class\_rank} can produce one-frame ID flickers
when two same-class bags are near each other and briefly swap x-order.
\texttt{class\_smooth} suppresses this by maintaining a vote history.

\subsubsection*{Majority-vote smoothing}
Each (camera, local\_id) pair maintains a \emph{vote history} of global ID
assignments:
\begin{equation}
  H_t^{(c,\ell)} = \bigl[g_{t-W+1},\; g_{t-W+2},\; \ldots,\; g_t\bigr],
  \quad W = 8.
\end{equation}
After \texttt{class\_rank} produces a raw assignment $g_t$, the smoothed
global ID is
\begin{equation}
  \boxed{\hat{g}_t^{(c,\ell)} = \text{mode}\!\left(H_t^{(c,\ell)}\right).}
  \label{eq:smooth}
\end{equation}
The window slides forward each frame; old votes are discarded.

\textbf{Effect}: a transient rank swap for $m < W/2$ frames does not change
the assigned global ID; only a swap that persists for more than half the
window actually propagates.

% =============================================================================
\section{Comparison of Online Methods}

\begin{table}[h]
\centering
\renewcommand{\arraystretch}{1.4}
\begin{tabularx}{\linewidth}{L{2.8cm} L{3.2cm} C{1.6cm} L{5.4cm}}
\toprule
\textbf{Method} & \textbf{Cost function} & \textbf{Latency} & \textbf{Failure mode} \\
\midrule
\code{class\_rank}
  & Ordinal x-rank (exact)
  & $O(N\log N)$
  & Same-class bags swap x-order → ID switch \\
\code{class\_iou}
  & $|\Delta\tilde{x}|$, Hungarian
  & $O(N^3)$ (LAPJV $O(N^2)$)
  & Sparse class → large $\Delta\tilde{x}$ → false match \\
\code{class\_smooth}
  & Ordinal x-rank + majority vote
  & $O(NW)$
  & Persistent order swap ($> W/2$ frames) \\
\bottomrule
\end{tabularx}
\caption{Online MCMOT associator comparison.  $N$: detections per class per
  frame; $W=8$: smoothing window.}
\end{table}

\begin{table}[h]
\centering
\renewcommand{\arraystretch}{1.4}
\begin{tabularx}{\linewidth}{L{2.4cm} C{1.3cm} C{1.3cm} C{1.3cm} C{1.2cm} L{4.8cm}}
\toprule
\textbf{Method} & \textbf{MOTA} & \textbf{IDF1} & \textbf{IDSW} & \textbf{fps} & \textbf{Note}\\
\midrule
\multicolumn{6}{l}{\textit{BoT-SORT / val set}} \\
\code{class\_rank}    & 96.2\% & 86.5\% & 31 & 0.9 & Fast; rank swaps cause IDSW \\
\code{class\_iou}     & 96.5\% & 89.5\% &  6 & 0.8 & Fewest IDSW; best IDF1 \\
\code{class\_smooth}  & 96.4\% & 86.8\% & 18 & 0.4 & Intermediate \\
\midrule
\multicolumn{6}{l}{\textit{BoT-SORT / test set}} \\
\code{class\_rank}    & 95.3\% & 83.5\% & 19 & 0.6 & \\
\code{class\_iou}     & 95.5\% & 83.9\% & 10 & 0.9 & Best IDF1 \\
\code{class\_smooth}  & 95.5\% & 83.6\% & 10 & 0.9 & \\
\bottomrule
\end{tabularx}
\caption{Online MCMOT results (MacBook M2, imgsz=1920).
  fps = scene-frames/s across all 4 cameras.  \code{class\_iou} consistently
  achieves the lowest ID-switch count and highest IDF1.}
\end{table}

% =============================================================================
\section{Offline vs.\ Online: Design Trade-offs}

\begin{table}[h]
\centering
\renewcommand{\arraystretch}{1.4}
\begin{tabularx}{\linewidth}{L{3.0cm} L{5.2cm} L{5.2cm}}
\toprule
& \textbf{Offline} & \textbf{Online} \\
\midrule
\textbf{Information}
  & Uses full tracklet history (past + future frames)
  & Only current active tracks \\
\textbf{Global optimality}
  & Can minimise cost over entire sequence
  & Greedy frame-by-frame \\
\textbf{Latency}
  & Post-processing only; not real-time
  & Real-time; suitable for live systems \\
\textbf{Best method}
  & \code{trk\_temporal} (IDF1 up to 98\%)
  & \code{class\_iou} (IDF1 up to 89.5\%) \\
\textbf{IDSW (best)}
  & 0 (\code{botsort}/test/\code{trk\_temporal})
  & 6 (\code{botsort}/val/\code{class\_iou}) \\
\bottomrule
\end{tabularx}
\caption{Summary comparison of offline vs.\ online MCMOT approaches.}
\end{table}

% =============================================================================
\section{References}

\begin{enumerate}[label={[\arabic*]},leftmargin=*,itemsep=2pt]
\item Y.\ Zhang, P.\ Sun, Y.\ Jiang, D.\ Yu, F.\ Weng, Z.\ Yuan, P.\ Luo,
      W.\ Liu, X.\ Wang.
      ``ByteTrack: Multi-Object Tracking by Associating Every Detection Box.''
      \textit{ECCV}, 2022.

\item N.\ Aharon, R.\ Orfaig, B.-Z.\ Bobrovsky.
      ``BoT-SORT: Robust Associations Multi-Pedestrian Tracking.''
      \textit{arXiv:2206.14651}, 2022.

\item F.\ Fleuret, J.\ Berclaz, R.\ Lengagne, P.\ Fua.
      ``Multicamera People Tracking with a Probabilistic Occupancy Map.''
      \textit{IEEE TPAMI}, 30(2):267--282, 2008.

\item Y.\ Bar-Shalom, T.\ E.\ Fortmann.
      \textit{Tracking and Data Association.}
      Academic Press, 1988.

\item S.\ Tang, B.\ Andres, M.\ Andriluka, B.\ Schiele.
      ``Multi-Person Tracking by Multicut and Deep Matching.''
      \textit{ECCV Workshops}, 2016.

\item A.\ Bewley, Z.\ Ge, L.\ Ott, F.\ Ramos, B.\ Upcroft.
      ``Simple Online and Realtime Tracking.''
      \textit{ICIP}, 2016.

\item H.\ W.\ Kuhn.
      ``The Hungarian Method for the Assignment Problem.''
      \textit{Naval Research Logistics Quarterly}, 2(1--2):83--97, 1955.
\end{enumerate}

\end{document}
